% Study 3, the question. Input from paper.tex; every number traces to sources.md and is asserted by
% check_numbers.py. Macros are provided in case the file is compiled before paper.tex defines them.
\providecommand{\qs}{QScore9}
\providecommand{\hc}{\texttt{hurdle\_count\_28}}

\section{Study 3: the question}
\label{sec:s3question}

Study 3 was registered on 2026-10-16 as item 7 of the protocol's dated additions, at commit
\texttt{3357d779db28cffa4048e7e333fd02b9434b4b6d} (full protocol SHA-256 with the addition
\texttt{05c7617b0391cbcca7bccddb6ff54a763d47ba88e5d31dac4bd98621a9721b29}), appended on top of commit
\texttt{04e92c1} after Study 2's results had been seen and before any Study 3 target value was accessed;
the only Study 3 source material in hand was the NYC 311 cell inventory of 2024, captured on 2026-09-14,
which holds no target-period count. Its hypothesis H3 is stated in the addition: on series whose context
zero fraction is at least 0.5, the substitution with a dynamic occurrence lowers \qs{} relative to the
model's own completed forecast; on series below 0.5 it raises it. Each stratum is tested on its own;
nothing is pooled across strata for confirmation; null, adverse, mixed and imprecise results receive the
same reporting. Where the addition says ``as registered'', Study 2's registration at \texttt{df55867}
applies.

The addition states why the study exists. Study 2 found, in its declared secondary stratum tables, that
replacing a frozen model's own zero mass with a context-only occurrence probability lowered \qs{} on
sparse npm series and raised it on dense ones, while the pooled confirmatory families spoke for the dense
regime, and a post-hoc analysis found that a dynamic occurrence estimate shrank the loss
(Section~\ref{sec:s2results}); that reading was descriptive. Study 3 tests it as a confirmatory family on
a second domain, with four changes from Study 2's construction, each answering a question that the
standing check of 2026-10-16, recorded in the repository as the referee's read, put to the paper. The tests
are per stratum, so that the confirmatory answer speaks for the regime it is asked about and not for
whatever dominates a pool (its item W1). The occurrence probability is the registered
\texttt{dynamic\_count} law's discounted estimate rather than the fixed 28-day Jeffreys estimate, which
ignores persistence (W9). The fourth cell of the occurrence-by-size factorial, the model's own occurrence
with a count positive law, is run as a declared secondary, since without it the two readings of the
accounting trade in Section~\ref{sec:discussion} cannot be separated (W2). And the primary count comparator
is the empirical hurdle, whose positive law can fit a dense series, so that the second contrast asks
whether the model's positive law keeps its value against a count law that can win, the question raised when
the registered dynamic count beat every foundation model on the grid (W4). What Study 3 does not answer is
fixed by the same registration: it tests no Markov occurrence estimate, adds no overdispersed cell to the
grid, says nothing about Wikipedia pageviews, runs the fourth cell on NYC 311 alone, and leaves the npm very
sparse stratum, with its nine series, as Study 2 left it.

The domain is NYC 311 service requests, as registered for Study 2's deferred slot: requests by creation day
for a problem type in a community board, a count of complaints about city services and so a third kind of
daily count beside GitHub stars and npm downloads. The claim the paper makes, stated in full in
Section~\ref{sec:discussion}, is the registered combination on the three settings it was run on, and
nothing broader.
